\documentclass[11pt]{article}
\usepackage[utf8]{inputenc}
\usepackage[T1]{fontenc}
\usepackage{amsmath,amssymb}
\usepackage[margin=1.1in]{geometry}
\usepackage{booktabs}
\usepackage{authblk}
\usepackage[colorlinks=true,linkcolor=blue,citecolor=blue,urlcolor=blue]{hyperref}

\title{Universality in Quantum Mechanics:\\
\large A Claim That Cannot Be Tested, and One That Can}

\author[1]{Claes Johnson}
\affil[1]{KTH Royal Institute of Technology, SE-100 44 Stockholm, Sweden \\
\small\texttt{claesjohnson@gmail.com}}

\date{\today}

\begin{document}
\maketitle

\begin{abstract}
Standard quantum mechanics is held to be universal for chemistry: one parameter-free equation,
with the nuclear charges as its only input, is taken to determine every atom, molecule and
material. We ask on what grounds that claim rests, and find that it cannot be tested. The
equation is a problem in $3N$ dimensions whose cost grows exponentially in the electron number;
beyond a few tens of electrons it is uncomputable, and what is computed in its place is not the
equation but a reduction carrying parameters --- basis sets, truncation orders, exchange--%
correlation functionals. Agreement between a parameterised reduction and observation does not
transfer to the parameter-free equation behind it, and the hierarchy of approximations cannot
launder it upward, since the hierarchy's exact limit is precisely the uncomputable object. The
universality of standard quantum mechanics is therefore verified at one to three electrons and
\emph{assumed} everywhere else, with the assumption protected from refutation by the
intractability of the thing assumed. We then ask the same question of RealQM, a model of
non-overlapping unit electron charge densities in three space dimensions, also parameter-free
with the nuclear charge as its only input. Because RealQM is computable at the scale of real
systems, its universality is an empirical claim: it can be put to atoms, molecules and solids and
compared with measurement without an intervening reduction. We report what that comparison gives
--- hydrogen exactly, helium to $0.24\%$, the first eighteen elements to $1.6\%$ in the mean, the
fourth row to $3.3\%$ --- and, with equal weight, where it fails: there is no spin, no dispersion,
and ionization energies are out of range. These are refutations of a testable claim, which is
what distinguishes it from the other. The asymmetry is dimensional, not technical: a model in
three dimensions can be held to account, a model in $3N$ cannot.
\end{abstract}

\section{Two parameter-free models}

Two models of electronic matter are on the table, and it is worth stating at the outset that both
are parameter-free in the same sense: the nuclear charge $Z$ is supplied and nothing is fitted.

\paragraph{Standard quantum mechanics.} The $N$-electron Schr\"odinger equation. The state is a
complex antisymmetric function $\Psi(x_1,\dots,x_N)$ on $\mathbb{R}^{3N}$, and the Hamiltonian is
kinetic energy plus Coulomb interaction among electrons and nuclei. No constant appears that was
not already in the description of the charges.

\paragraph{RealQM.} The state is a collection of $N$ one-electron charge densities on
non-overlapping domains in $\mathbb{R}^3$, continuous across free interfaces, interacting by
Coulomb potentials without self-repulsion, with the ground state obtained by minimising the
energy. Again no constant appears beyond $Z$.

The two differ in the dimension of the object they carry: $3N$ against $3$. Everything that
follows turns on that difference, and the burden of this article is that it is not a difference
of convenience.

Both models are advanced as \emph{universal}, so we had better say what the claim asserts.

\section{What a universality claim asserts}\label{sec:what}

A model is universal over a domain of phenomena if three conditions hold.

\begin{enumerate}
\item \textbf{Species-blindness.} The model contains no reference to the kind of system it is
applied to. It does not mention atoms, bonds, metals or proteins. One writes down the same thing
for helium and for a protein, differing only in $Z$ and the geometry.
\item \textbf{Closure under composition.} Joining two systems yields the same model for the union.
There is no boundary in the domain at which one must change theories.
\item \textbf{Accountability.} The model's predictions can be compared with measurement across
the domain, and the comparison can go against it.
\end{enumerate}

The first two are formal and are easily checked. The third is not formal at all, and it is the one
that does the work: a model satisfying only the first two is a \emph{framework}, not a universal
theory, because nothing it says about most of its domain can be held to anything.

Standard quantum mechanics satisfies (1) and (2) perfectly. Our claim is that over almost all of
its asserted domain it does not satisfy (3), and cannot be made to.

\section{Why the universality of standard quantum mechanics cannot be tested}

\subsection{The equation is uncomputable where the claim is interesting}

The state is a function on $\mathbb{R}^{3N}$. Representing it on a grid or in a basis costs
exponentially in $N$: full configuration interaction in a basis of $M$ functions scales as
$\binom{M}{N}^2$ in the worst case, and in practice the exact solution is out of reach past
roughly twenty electrons. This is not a matter of waiting for better machines. The ground-state
problem for a local Hamiltonian is QMA-complete, so no efficient general algorithm is believed to
exist even on a quantum computer. The barrier has the character of a theorem, not of an
engineering limit.

The systems for which universality is an interesting claim --- a catalyst, a protein, a crystal,
a cell --- have $10^3$ to $10^{23}$ electrons. For every one of them the model cannot be solved,
now or in principle by any method we expect to exist.

\subsection{What is computed instead is not the model}

Since the equation cannot be solved, it is not the equation that is compared with experiment.
What is compared is a reduction, and every reduction in use carries parameters.

\begin{itemize}
\item \textbf{Basis sets.} Gaussian basis sets are constructed by optimising exponents and
contraction coefficients, atom by atom, against atomic calculations. The complete-basis-set limit
is reached by extrapolation through a fitted functional form.
\item \textbf{Truncation.} Coupled-cluster theory truncates the excitation operator; CCSD(T)'s
perturbative triples correction is not variational and has no bound. Its celebrated accuracy of
about $1$~kcal/mol is an empirical regularity over well-behaved closed-shell molecules, and it
degrades without warning for multireference systems.
\item \textbf{Functionals.} Density functional theory, which performs essentially all practical
calculation on large systems, chooses an exchange--correlation functional from a zoo. B3LYP, PBE,
M06 and $\omega$B97X are different theories with different failures, and several are fitted to
benchmark sets outright.
\end{itemize}

We should be fair about the strongest case. There is a window --- small closed-shell molecules of
up to a few tens of atoms --- in which CCSD(T) near the basis-set limit is nearly parameter-free
and agrees with measurement to about a kcal/mol across a wide range of chemistry. That is a real
and impressive achievement and we do not dispute it. But it is a window, and outside it, which
is to say for almost everything, the computed object is a functional-bearing reduction.

\subsection{The hierarchy cannot transfer verification upward}

It is often said that this does not matter, because the approximations form a convergent
hierarchy --- Hartree--Fock, MP2, CCSD, CCSD(T), full CI --- whose limit is the exact solution.
The hierarchy therefore appears to connect what is computed to what is claimed.

It does not, for a reason that is structural. The hierarchy has two ends and both are unavailable.
Its exact limit is the uncomputable object itself, so one can never stand at the top and check.
Its lower rungs, which are what gets computed, each carry parameters, so agreement there is
agreement of a parameterised procedure. There is no rung which is both reachable and
parameter-free, and consequently no step at which empirical success attaches to the
parameter-free equation rather than to the reduction.

The existence of a convergent sequence establishes that the reductions \emph{approach} something.
It does not establish that the something is nature, because that identification is exactly what
was to be tested.

\subsection{The practice this produces is unfalsifiable}

The consequence is visible in how discrepancies are handled. When DFT fails on dispersion ---
and semilocal and hybrid functionals do fail on it, systematically --- the field's response is to
add Grimme's empirical $D3$ or $D4$ correction, a fitted term. The failure is charged to the
functional. When a coupled-cluster result is wrong for a transition-metal complex, it is charged
to the truncation or the basis. In no case is it charged to the Schr\"odinger equation, and this is
not dogmatism: it is forced, because no computation of the equation exists against which the
equation could be found wanting.

A claim that cannot be contradicted by any measurement is not thereby strong. It has been removed
from the class of claims that measurement bears on. Over the great majority of its asserted
domain, that is the position of the universality of standard quantum mechanics: verified
superbly at one to three electrons, and assumed beyond, with the assumption shielded by the
intractability of the object assumed.

\section{Why the universality of RealQM can be tested}

RealQM is in the opposite position on exactly the point that matters, and for exactly one reason:
it is computable where the claim is interesting.

The state is $N$ densities on $\mathbb{R}^3$. The cost is the number of electrons times the number
of mesh points --- linear in $N$, not exponential. A molecule, a cluster, a solid is the same
computation with more kernels and more domains. So the model itself, with $Z$ as its only input
and no reduction interposed, can be run on the systems in question and the output compared with
measurement.

That makes universality an empirical question about RealQM, in the sense of
Section~\ref{sec:what}(3). The claim is accountable. What follows is the account.

\subsection{What it gets right}

\paragraph{Hydrogen is an identity, not a test.} For one electron there is no exclusion and no
partition, and RealQM \emph{is} the Schr\"odinger equation. Both give $E=-\tfrac12$ and
$\alpha = 4.5\,a_0^3$. Nothing is learned about either model here, and we record it only because
it locates where the two can first differ: at two electrons, where the exclusion postulate enters.

\paragraph{Helium.} Two electrons is the first real test, and the exact answer is known.

\begin{center}
\begin{tabular}{lrr}
\toprule
 & $E$ (Ha) & error \\
\midrule
exact (nonrelativistic)                 & $-2.9037$ & --- \\
RealQM, two domains, front-tracked plane & $-2.91$   & $0.24\%$ \\
Hartree--Fock                            & $-2.8617$ & $1.4\%$ \\
RealQM, fixed Neumann plane              & $-2.99$   & $3\%$ \\
\bottomrule
\end{tabular}
\end{center}

The front-tracked two-domain result is closer to exact than Hartree--Fock is. We draw no triumph
from this --- the comparison is one number --- but it does dispose of the expectation that
replacing antisymmetry by territorial exclusion must be grossly wrong. On helium's energy it is
not.

\paragraph{The first eighteen elements, and the fourth row.} Computed all-electron on the bare
kernel $-Z/r$, with no core radius, no fitted length and nothing adjusted per element: the mean
relative error is $1.6\%$ with maximum $4.5\%$ (silicon), and $13$ of $17$ atoms fall within
$2.5\%$. Extended to the fourth row at a mesh chosen by $Z$, the mean is $3.3\%$ with maximum
$6.7\%$ (selenium), and all eighteen entries are under-bound --- the direction the spherical
reduction predicts. These are the errors of those meshes and not a converged accuracy; single
entries move by a few per cent under refinement, systematically rather than randomly, so the mean
should be read at the same level as the entries.

\paragraph{The periodic table's capacities.} Two filled packings appear and no third: the
antipodal pair, $2$, and the octet, $8=4{+}4$, two dual tetrahedra. Row capacities are sums of
these alone --- $2$, $8$, $18=2{+}8{+}8$, $32=8{+}8{+}8{+}8$ --- and a period ends when its
capacity is exhausted, giving the closures $2,10,18,36,54,86,118$.

\subsection{What it gets wrong, stated with the same weight}

A testable claim earns its standing by the failures it records, so these belong in the same
section and not in a footnote.

\begin{itemize}
\item \textbf{There is no spin.} The singlet--triplet splitting of helium, ortho and para states,
magnetism, Hund's rules: these are not approximated badly, they are absent. No refinement reaches
them.
\item \textbf{There is no dispersion.} Two neutral closed-shell atoms should attract as $R^{-6}$.
Dispersion is the energy of \emph{correlated fluctuations}, and RealQM's densities are static, so
it should have none. We note that Hartree--Fock has none either --- it gives He$_2$ purely
repulsive --- and that practical DFT supplies it by a fitted term. The He$_2$ well is
$0.035$~mHa at $R=5.6\,a_0$, and measuring it is the sharpest available test of whether the
free interfaces contain any correlated motion at all.
\item \textbf{Ionization energies are out of range.} A total energy is a bulk sum and is robust; an
ionization energy depends on a single loosely bound charge and is fragile. Proton affinities come
out with the wrong sign.
\item \textbf{Polarisability is unresolved.} It is the one quantity here checked against
spectroscopy rather than against another calculation, and we do not yet have a converged value
for helium: the available solvers disagree, and the number previously reported rested on a
configuration that carried no electron--electron repulsion. We state this rather than quote it.
\item \textbf{There is no convergence ladder.} RealQM is a fixed ansatz. Refining its mesh
converges to the ansatz's own answer, not toward nature's. There is no next order to add.
\end{itemize}

The last of these is the deepest, and it bounds the claim: RealQM can assert universal
\emph{scope} but not universal \emph{fidelity}. It applies everywhere, with an error floor of a
few per cent and no route down.

\section{The asymmetry is dimensional}

It would be easy to read the foregoing as a comparison of accuracy, and that is not the point.
The point is that the two universality claims are of different \emph{kinds}, and the dimension of
the state is what makes them so.

A model whose state lives in $3N$ dimensions cannot be evaluated on a system of chemical
interest. Its universality is therefore a statement about an object that is never exhibited: an
existence claim without a witness. A model whose state lives in $3$ dimensions can be evaluated,
so its universality is a prediction, and predictions can fail.

This is the same distinction that separates a non-constructive existence proof from a computed
solution, and it recurs across applied mathematics. The symbolic object is more general and
cannot be inspected; the computed object is specific and can be checked. Neither is worthless,
but they are not interchangeable, and a claim of universality means something different in each.

Set out plainly, each model holds the strong half of a different trade, and neither half is
redeemed:

\begin{center}
\begin{tabular}{lll}
\toprule
 & standard QM & RealQM \\
\midrule
state                      & $3N$ dimensions        & $3$ dimensions \\
parameters                 & none (in the equation) & none \\
cost                       & exponential in $N$     & linear in $N$ \\
species-blind              & yes                    & yes \\
closed under composition   & yes                    & yes \\
verified range             & $1$--$3$ electrons     & $1$--$36$ electrons \\
verified precision         & few parts in $10^5$    & few per cent \\
exact in principle         & yes                    & no \\
convergence ladder         & yes, but unreachable   & none \\
\textbf{universality}      & \textbf{assumed}       & \textbf{tested} \\
\bottomrule
\end{tabular}
\end{center}

Standard quantum mechanics is verified at three electrons and asserted at $10^{23}$. RealQM is
verified at thirty-six and asserted beyond. Over the range where either is actually used, the
second assertion is the better evidenced of the two --- which is a remark about evidence, not
about truth.

\section{What this does and does not establish}

It does not establish that standard quantum mechanics is false. Where it can be computed it is
superb, and nothing in this article gives any reason to doubt it. It does not establish that
RealQM is true; RealQM is demonstrably incomplete, in the ways listed above, and its lack of a
ladder means its errors cannot be driven down.

What it establishes is narrower and, we think, more useful. The conventional ranking of these two
models treats the $3N$ formulation as the serious one and the $3$-dimensional model as an
approximation to be tolerated. On the specific question of universality that ranking inverts,
because universality is a claim about scope and accountability, and only one of the two models is
accountable over the scope it claims. A framework that applies in principle to everything and can
be checked against nothing has a weaker title to universality than a model which applies to
everything, is checked, and is sometimes wrong.

\section{The test to run}

If the argument above is right, then the productive question is not whether RealQM is universal
in the abstract but which of its failures are structural and which are merely unfinished. One of
them is decisive and cheap.

Dispersion is the natural place, because it is where a static model should fail and where the
dominant practical method also fails, patching the gap with a fitted term. If RealQM's free
interfaces --- which are, unlike orbitals, real movable surfaces --- were allowed to fluctuate in
a correlated way, that is precisely what a second rung on the ladder would look like, and
precisely what dispersion is. So He$_2$ is not merely a missing number. It is the test of whether
RealQM has a second rung at all, and therefore of whether its universality can ever be a claim
about fidelity and not only about scope.

\end{document}
